% SEGMENT OLP-0545-S01
% Part: set-theory
% Chapter: z
% Section: nat

\documentclass[../../../include/open-logic-section]{subfiles}

\begin{document}

\olfileid{sth}{z}{nat}
\olsection{நமது இயல் எண்களைத் தேர்ந்தெடுத்தல்}

%முறைசாராக் கொள்கையான \stagesinf{} ஐ நியாயப்படுத்தலாம் என எடுத்துக்கொள்வோம்.
%அந்த முறைசாராக் கொள்கையிலிருந்து நாம் பெற்ற குறிப்பிட்ட அடிகோளைப் பற்றிக்
%கருத்துரைப்பதும் பயனுள்ளது.

\olref[infinity-again]{defnomega} இல், ``இயல் எண்கள்'' என்ற தொடரை
வெளிப்படையாக வரையறுத்தோம். இந்த நிர்ணயத்தை நீங்கள் எவ்வாறு புரிந்துகொள்ள
வேண்டும்? இது மெய்ப்பொருளியக் கூற்று அல்ல; சில கணங்களை இயல் எண்களாக
\emph{கருதுவதற்கான} ஒரு முடிவு மட்டுமே. அவ்வாறு கருதுவதற்கான காரணங்களை
\olref[sfr][rel][ref]{sec}, \olref[sfr][arith][ref]{sec},
\olref[sfr][infinite][dedekindsproof]{sec} ஆகியவற்றில் தொட்டுக்காட்டினோம்.
ஆனால் இக்காரணங்களை இன்னும் கூர்மையாக்கலாம்.

நமது முடிவிலி அடிகோள் \citet{VonNeumann1925} ஐப் பின்பற்றுகிறது. ஆனால்
அதற்குப் பதிலாக நாம் ஏற்றிருக்கக்கூடிய வேறோர் அடிகோள் இதோ:

\begin{defish}
\emph{செர்மேலோவின் \citeyear{Zermelo1908Untersuchungen} ஆம் ஆண்டு
முடிவிலி அடிகோள்.} $\emptyset \in A$ ஆகவும்
$(\forall x \in A)\{x\} \in A$ ஆகவும் உள்ள ஒரு கணம் $A$ உள்ளது.
\end{defish}

நமது (ஃபான் நாய்மன் வழியமைந்த) முடிவிலி அடிகோளுக்குப் பதிலாக
செர்மேலோவின் அடிகோளைப் பயன்படுத்தியிருந்தாலும், அதேபோல் ஒரு
டெடெகிண்ட்-முடிவிலிக் கணமும் அதனால் ஒரு டெடெகிண்ட் இயற்கணிதமும் நமக்குக்
கிடைத்திருக்கும். செர்மேலோவின் அணுகுமுறையில், நமது இயற்கணிதத்தின்
சிறப்பிக்கப்பட்ட உறுப்பு மீண்டும் $\emptyset$ ஆகவே (அதாவது, $0$ க்கான
நமது பதிலியாக) இருந்திருக்கும்; ஆனால் ஒன்றுக்கொன்றான சார்பு
$x \mapsto x \cup \{x\}$ என்பதற்குப் பதிலாக $x \mapsto \{x\}$ என்ற
வரைப்படத்தால் தரப்பட்டிருக்கும். இதன் மிக எளிய விளைவு: செர்மேலோ
$2$ ஐ $\{\{\emptyset\}\}$ எனக் கருதுகிறார்; ஆனால் (ஃபான் நாய்மனைப்
பின்பற்றும்) நாம் $2$ ஐ $\{\emptyset, \{\emptyset\}\}$ எனக் கருதுகிறோம்.

ஒரு முடிவிலி அடிகோளை விட மற்றொன்றை ஏன் தேர்ந்தெடுக்க வேண்டும்? ஃபான்
நாய்மனின் அணுகுமுறை முடிவிலிக்கடந்த எண்களைக் கையாள நன்கு ``விரிவடைகிறது''
என்பதே முதன்மையான நடைமுறைக் காரணம். \olref[ordinals][]{chap} தொடங்கி
இதனை ஆராய்வோம். எனினும், \emph{எண்கணிதம் செய்வது} என்ற எளிய
கண்ணோட்டத்தில், இரு அணுகுமுறைகளும் ஒரே அளவு நன்கு செயல்படும். ஆகவே இயல்
எண்கள் கணங்களாக \emph{உள்ளன} என்று ஒருவர் கூறினால், வெளிப்படையான கேள்வி:
\emph{அவை எந்தக் கணங்கள்?}

இந்தத் துல்லியமான கேள்வியை \citet{Benacerraf1965} புகழ்பெறச் செய்தார்.
ஆனால் நாம் ஏற்கெனவே பலமுறை எதிர்கொண்ட ஒரு நிகழ்வின் மிகப் புகழ்பெற்ற
எடுத்துக்காட்டு மட்டுமே இது என்பதை வலியுறுத்துவது பயனுள்ளது. அடிப்படைக்
கருத்து இதுதான். ``உள்ளுணர்வான'' பலவகைப் பொருள்களை \emph{போலமைக்க} கணக்
கோட்பாடு நமக்கு ஒரு வழி தருகிறது: மெய்யெண்கள், விகிதமுறு எண்கள், முழுக்கள்,
இயல் எண்கள் மட்டுமல்லாமல், வரிசைப்படுத்தப்பட்ட ஜோடிகள், சார்புகள்,
தொடர்புகளையும் போலமைக்கிறது. ஆனால் கணக் கோட்பாடு ஒருபோதும்
\emph{தனித்ததொரு} போலமைப்புத் தேர்வை நமக்குத் தருவதில்லை. அதே அளவு
நன்கு பயன்பட்டிருக்கக்கூடிய மாற்றுகள் \emph{எப்போதும்} உள்ளன.
% SOURCE CORRECTION TA-STH-008: repair the source's “this significance” to “the significance”.

\end{document}
